\part{Percolation / Random-Cluster Model}

\section{Probability on lattices}

\begin{exercise}[Independent edges, monotone events, and lattice probability]
Let $\mathbb E^d$ be the nearest-neighbour edges of $\Z^d$.
A bond configuration is $\omega\in\{0,1\}^{\mathbb E^d}$;
an edge is open when its value is one. The measure $\Prob_p$
makes the edge values independent with probability $p$ of
being open. An event is increasing if opening more edges
cannot destroy it. A cylinder event depends on finitely many
edges.
\begin{enumerate}[label=(\alph*)]
\item Construct the monotone coupling
$\omega_e(p)=\mathbf1_{U_e\le p}$ from independent uniform
variables. Prove that probabilities of increasing events
are nondecreasing in $p$.
\item Prove the Harris inequality
$\Prob_p(A\cap B)\ge\Prob_p(A)\Prob_p(B)$ for increasing
cylinder events, by induction on the number of edges and
conditioning on the last edge. Extend it to increasing
connection events by monotone limits.
\item An edge is pivotal for a cylinder event when changing
that edge alone changes whether the event occurs.
Differentiate first with independent densities $p_e$ and
then set $p_e=p$ to prove Russo's formula
$\frac{\dd}{\dd p}\Prob_p(A)=\sum_e\Prob_p(e\text{ pivotal})$
for increasing cylinder events.
\end{enumerate}
\end{exercise}

\begin{exercise}[Random-cluster weights and the Edwards--Sokal coupling]
On a finite graph $(V,E)$ write $o(\omega)$ for the number
of open edges and $k(\omega)$ for the number of connected
components, including isolated vertices. For $Q>0$ define
\[
 \phi_{p,Q}(\omega)=Z_{\rm RC}^{-1}
    p^{o(\omega)}(1-p)^{|E|-o(\omega)}Q^{k(\omega)}.
\]
A boundary wiring is an equivalence relation on specified
boundary vertices; count components after making those
identifications. For integer $Q\ge2$ the Potts Hamiltonian
is $H=-J\sum_{ij}\mathbf1_{\sigma_i=\sigma_j}$,
$\sigma_i\in\{1,\dots,Q\}$.
\begin{enumerate}[label=(\alph*)]
\item Put $p=1-e^{-\beta J}$ and assign joint weight
$\prod_{ij}[(1-p)\mathbf1_{\omega_{ij}=0}
+p\mathbf1_{\omega_{ij}=1}\mathbf1_{\sigma_i=\sigma_j}]$.
Prove that its spin marginal is the Potts measure and its
edge marginal is $\phi_{p,Q}$, with
$Z_{\rm RC}=e^{-\beta J|E|}Z_{\rm Potts}$.
\item Conditional on edges, colour each cluster independently
and uniformly. Derive
\[
 \langle\mathbf1_{\sigma_x=\sigma_y}\rangle-\frac1Q
   =\left(1-\frac1Q\right)\phi_{p,Q}(x\leftrightarrow y).
\]
For $Q=2$ identify the Ising coupling $K=\beta J/2$
and prove $\langle\sigma_x\sigma_y\rangle
=\phi_{p,2}(x\leftrightarrow y)$.
\item For $Q\ge1$ prove supermodularity of $k$ by considering
the effect of opening two edges. Apply the FKG lattice
condition of the section Boundary conditions to prove
positive association. Calculate the conditional open-edge
probability as $p$ if its endpoints are already connected
and $p/(p+Q(1-p))$ otherwise.
\end{enumerate}
\end{exercise}

\section{Connectivity}

\begin{exercise}[Open clusters and bounds on the percolation threshold]
Let $C(x)$ be the cluster containing $x$, let
$\theta(p)=\Prob_p(|C(0)|=\infty)$, and put
$p_c=\sup\{p:\theta(p)=0\}$.
For $\Lambda_n=\{-n,\dots,n\}^d$, define its boundary as
in the section Thermodynamic limit.
\begin{enumerate}[label=(\alph*)]
\item Prove that $\{|C(0)|=\infty\}$ is the decreasing
intersection of $\{0\leftrightarrow\partial\Lambda_n\}$.
Show that an infinite cluster contains an infinite
self-avoiding path by choosing successive extensions in
the finitely branching tree of finite paths.
\item Count at most $2d(2d-1)^{n-1}$ self-avoiding
$n$-edge paths from zero. Deduce
$p_c(\Z^d)\ge1/(2d-1)$ for $d\ge2$, and prove
$p_c(\Z)=1$ separately.
\item On $\Z^2$, if the origin's cluster is finite, trace a
closed dual circuit of closed primal edges around it.
Bound the number of possible length-$n$ circuits by
$4(2n+1)^2\,3^{n-1}$, by specifying a nearby starting edge
and a nonbacktracking continuation. Deduce $\theta(p)>0$
when $p$ is sufficiently close to one and hence $p_c<1$.
\end{enumerate}
\end{exercise}

\begin{exercise}[Planar duality and random-cluster self-duality]
Let $G$ be a connected finite plane graph and let $G^*$
have one vertex for each face, including the unbounded face.
Declare the dual edge open precisely when its crossed
primal edge is closed.
\begin{enumerate}[label=(\alph*)]
\item Apply Euler's formula to the open subgraph to prove
$k(\omega)=k(\omega^*)+o(\omega^*)+|V|-|E|-1$.
\item Substitute this equality in the random-cluster weight
and derive
\[
 \frac{p^*}{1-p^*}=\frac{Q(1-p)}p,\qquad
 p_{\rm sd}=\frac{\sqrt Q}{1+\sqrt Q}.
\]
Explain the exchange of a free boundary with a single wired
outer dual vertex through the graph construction.
\item For a rectangle prove that exactly one of a primal
left-right open crossing and a dual top-bottom open
crossing occurs, with the dual boundary arcs specified
by the crossed boundary edges. Use planar separation.
Distinguish this finite combinatorial identity from an
assertion that $p_{\rm sd}=p_c$ in an infinite lattice.
\end{enumerate}
\end{exercise}

\section{Critical phenomena}

\begin{exercise}[An exact percolation transition on a branching lattice]
Consider independent bond percolation on the rooted tree
with $b\ge2$ children at each vertex. Write $\theta_b(p)$
for the probability that the root has infinitely many
descendants in its open cluster.
\begin{enumerate}[label=(\alph*)]
\item Let $x_n$ be the probability of an open path from
the root to generation $n$. Prove $x_0=1$,
$x_{n+1}=1-(1-px_n)^b$, and $x_n\downarrow\theta_b(p)$.
\item Prove that $\theta_b(p)$ is the largest solution in
$[0,1]$ of $x=1-(1-px)^b$. Use strict concavity to show
$p_c=1/b$, with zero survival at $p_c$.
\item Expand at the bifurcation to obtain
$\theta_b(p)\sim 2b^2(p-1/b)/(b-1)$ as $p\downarrow1/b$.
Count open descendants by generations to prove
$\E_p|C(0)|=(1-bp)^{-1}$ for $p<1/b$.
\end{enumerate}
\end{exercise}

\section{Correlation length}

\begin{exercise}[Exponential connectivity decay and the inverse correlation seminorm]
For $0<p<1$ on $\Z^d$ set
$\tau_n(x)=-\log\Prob_p(0\leftrightarrow nx)$.
Use the correlation-length convention of the section
Correlation functions, with spin correlation replaced by
connection probability.
\begin{enumerate}[label=(\alph*)]
\item Use Harris' inequality and translation invariance
to prove $\tau_{n+m}(x)\le\tau_n(x)+\tau_m(x)$.
Prove the subadditive limit formula
$\tau(x)=\lim_n\tau_n(x)/n=\inf_n\tau_n(x)/n$ by division
with remainder.
\item Prove integer homogeneity, the triangle inequality,
symmetry, and $\tau(x)\le-\log p\,|x|_1$. Extend to a
continuous seminorm on $\R^d$.
If $\xi^{-1}=\tau(e_1)$, use coordinate-reflection symmetry
to prove $\tau(x)\ge|x|_\infty/\xi$ and hence
$\Prob_p(0\leftrightarrow x)\le e^{-|x|_\infty/\xi}$.
\item Count self-avoiding paths to prove $\xi<\infty$
for $p<1/(2d-1)$. With
$\chi(p)=\E_p|C(0)|=\sum_x\Prob_p(0\leftrightarrow x)$,
derive $\chi(p)\le C_d\max(1,\xi)^d$ whenever $\xi<\infty$.
\item On the rooted tree a fixed length-$n$ path has
connection probability $p^n$. Compute its decay length
$-1/\log p$ and show it stays finite at $p=1/b$ even though
the mean cluster size diverges. Identify the exponential
growth in the number of sites per generation.
\end{enumerate}
\end{exercise}

\section{Scaling and universality}

\begin{exercise}[Critical exponents and finite-size crossing windows]
When the indicated asymptotics hold, define exponents by
$\theta(p)\sim C(p-p_c)^{\beta_{\rm perc}}$,
$\chi(p)\sim C'(p_c-p)^{-\gamma}$, and
$\xi(p)\sim C''|p-p_c|^{-\nu}$ on the specified sides.
Universality of an exponent means its equality for a
specified class of models after identifying the observables;
it does not assert equality of thresholds or amplitudes.
\begin{enumerate}[label=(\alph*)]
\item Determine $\beta_{\rm perc}$ and $\gamma$ on the
rooted branching lattice from the section Critical phenomena.
Explain why no divergent path-decay exponent $\nu$ is
obtained there.
\item Let $\pi_n(p)$ be finite-volume crossing probabilities.
Suppose, as mathematical data for this conditional assertion,
$\pi_n(p_c+t n^{-1/\nu})\to\Phi(t)$, locally uniformly
in $t$, where $\Phi$ is continuous and strictly increasing
from zero to one. Put
$p_n(a)=\inf\{p:\pi_n(p)\ge a\}$.
Use monotonicity to prove
\[
 n^{1/\nu}(p_n(b)-p_n(a))
 \longrightarrow \Phi^{-1}(b)-\Phi^{-1}(a)
 \quad(0<a<b<1).
\]
\item For a change of parameter $\widetilde p=f(p)$ with
$f'(p_c)>0$, determine which exponents, amplitudes, thresholds,
and window coefficients change. State the normalization
needed when comparing crossing windows on two lattices.
\end{enumerate}
\end{exercise}

\section{SLE / conformal invariance}

\begin{exercise}[Loewner evolution and Brownian scaling]
For a given continuous real driving function $W$ define
$g_t(z)$ in the upper half-plane, up to its collision lifetime,
by
$\partial_tg_t(z)=2/(g_t(z)-W_t)$, $g_0(z)=z$.
A standard Brownian motion is a continuous process $B_0=0$
with independent Gaussian increments of mean zero and
variance $t-s$ on an interval $[s,t]$.
Given such a process, the random Loewner evolution driven
by $W_t=\sqrt\kappa B_t$ is chordal $\mathrm{SLE}_\kappa$.
\begin{enumerate}[label=(\alph*)]
\item Solve the evolution for $W=0$ and obtain
$g_t(z)=\sqrt{z^2+4t}$, with the branch asymptotic to $z$.
Determine the removed vertical slit.
\item For general continuous $W$ prove, by rescaling the ODE,
that $\widetilde g_t(z)=a g_{t/a^2}(z/a)$ is driven by
$\widetilde W_t=aW_{t/a^2}$.
Use Gaussian increment laws to prove scale invariance
of $\mathrm{SLE}_\kappa$ up to the corresponding time change.
\item For deterministic $s$, compose $g_{s+t}$ with
$g_s^{-1}(\,\cdot+W_s)$ and subtract $W_s$.
Show that the resulting evolution is driven by
$W_{s+t}-W_s$, independent of the past with the original
driving law. Prove injectivity and existence of the inverse
on the surviving domain using the forward and reverse ODEs.
\end{enumerate}
\end{exercise}

\begin{exercise}[Prediction: conformally invariant percolation crossing probabilities]
Let $D$ be a Jordan domain with four counterclockwise marked
boundary points $a,b,c,d$, and fix a conformal map to the
upper half-plane sending them to $0,x,1,\infty$, $0<x<1$.
Use independent site percolation at density $1/2$ on a
triangular lattice of mesh $\delta$ approximating $D$.
The observable is a black crossing from arc $(ab)$ to
arc $(cd)$. For this model assume, as proved by Smirnov
\cite{smirnov-percolation}, that its probability converges to
\[
 P(x)=\frac{\displaystyle\int_0^x
                  t^{-2/3}(1-t)^{-2/3}\dd t}
             {\displaystyle\int_0^1
                  t^{-2/3}(1-t)^{-2/3}\dd t}.
\]
\begin{enumerate}[label=(\alph*)]
\item Prove $P(x)+P(1-x)=1$, $P(1/2)=1/2$, and
$P(x)\sim3x^{1/3}/\int_0^1t^{-2/3}(1-t)^{-2/3}\dd t$
as $x\downarrow0$. Calculate $P(1/4)$ numerically.
\item Use a conformal symmetry exchanging opposite pairs
of sides to derive crossing probability $1/2$ for a square.
Specify that $x$ is the cross-ratio of the mapped real
boundary points, not of arbitrary points in the original domain.
\item For $M$ independent lattice samples compute the mean
and variance of the empirical crossing frequency.
Separate this sampling uncertainty from the finite-mesh
error, for which the stated hypothesis supplies no rate.
State precisely the triangular-site model covered by the
hypothesis; a corresponding square-bond assertion requires
an additional universality result.
\end{enumerate}
\end{exercise}
